\chapter{Evidencia complementaria}
\label{anexo:evidencia}

\section{Inventario de artefactos}

Los estudios persisten bastante más evidencia de la que el texto principal utiliza. Se documenta
aquí para auditoría y para facilitar replicaciones, con sus dimensiones reales, de modo que quien
quiera verificar una cifra sepa dónde buscarla.

Conviene recordar la regla de procedencia, porque afecta a dónde hay que mirar cada cifra: los
artefactos \emph{predictivos} pertenecen al Model Study 3 (\studyidC) y los \emph{económicos} del
documento final proceden del ganador del Portfolio Study (\studyportfolio), bajo
\artefacto{evidence\_best\_full/}. Las dimensiones que siguen corresponden al Model Study 3.

\begin{description}[leftmargin=0.5cm,style=nextline,font=\normalfont\ttfamily\small]
  \item[evidence/equity.parquet] 135 filas $\times$ 15 columnas. Serie mensual de la cartera y del
  \textit{benchmark} entre abril de 2015 y junio de 2026, con retornos, rotación, arrastre de costes
  y peso en efectivo por snapshot.
  \item[evidence/annual\_metrics.parquet] 12 $\times$ 9. Resultado por año, de 2015 a 2026.
  \item[evidence/orders.parquet] 464 $\times$ 12. Cada orden emitida, con su motivo declarado,
  precios, comisión, \textit{slippage} y resultado realizado.
  \item[evidence/positions.parquet] 982 $\times$ 12. Posiciones por snapshot, con precio de entrada,
  valoración, unidades, meses en cartera, percentil actual y rango del meta.
  \item[evidence/meta\_weights.parquet] 675 $\times$ 6. Peso asignado a cada agente en cada fecha,
  con el estado que indica si fue aprendido o de reparto uniforme.
  \item[evidence/rank\_ic\_diagnostics.parquet] 861 $\times$ 6. Rank-IC por agente y fecha: siete
  señales sobre 123 cohortes.
  \item[evidence/rank\_tail\_diagnostics.parquet] 123 $\times$ 16. Comportamiento de las colas y
  concentración de los pesos del meta-agente.
  \item[evidence/signal\_health.parquet] 135 $\times$ 7. Rank-IC contraído observable en cada
  snapshot, con memorias de ocho y dieciséis trimestres.
  \item[evidence/feature\_diagnostics.parquet] 73 $\times$ 6. Cobertura, Rank-IC univariante e
  importancia por variable.
  \item[evidence/agent\_scores.parquet] 60.380 $\times$ 26. Puntuación de cada agente para cada
  valor y snapshot.
  \item[evidence/model\_feature\_attribution.parquet] 1.613 $\times$ 6. Coeficientes por agente y
  fecha de reajuste.
  \item[evidence/agent\_local\_attribution.parquet] 1.309.618 $\times$ 10. Contribución de cada
  variable a la puntuación de cada acción en cada fecha, con dirección y orden de importancia. Es el
  artefacto más voluminoso del estudio y sostiene el análisis de explicabilidad de la
  Sección~\ref{sec:atribucion-local}. El manuscrito lo explota sólo de forma agregada: queda
  disponible el detalle por acción individual, que permitiría explicar decisiones concretas de
  cartera.
  \item[evidence/signal\_calibration.parquet] 60.380 $\times$ 6. Alfa esperado frente a alfa
  realizado, que permite comprobar si la calibración de percentil a rentabilidad está sesgada.
  \item[evidence/profiles/] Ocho perfiles, cada uno con su propia curva de patrimonio, métricas
  anuales, órdenes y posiciones.
  \item[robustness.json, attribution.json, decisions.json] Contrastes de robustez, atribución
  factorial y traza completa de decisiones con todos sus candidatos.
\end{description}

\subsection*{Artefactos del Portfolio Study}

\begin{description}[leftmargin=0.5cm,style=nextline,font=\normalfont\ttfamily\small]
  \item[portfolio\_grid.parquet] 1.728 $\times$ 19. Una fila por cartera evaluada, con sus seis
  coordenadas y sus métricas medidas sobre la serie recortada en 2024. De las descartadas se
  conserva únicamente esta fila de resumen.
  \item[portfolio\_winner.json] Combinación ganadora, su resumen de selección, su confirmación sobre
  la serie completa y la declaración explícita del año de corte de la rejilla.
  \item[portfolio\_profiles.parquet] Los ocho perfiles evaluados con la cartera ganadora, cada uno
  medido dos veces: sobre la serie recortada y sobre la completa.
  \item[evidence\_best\_full/] Evidencia económica completa del ganador sobre la serie entera ---
  patrimonio, órdenes, posiciones y métricas anuales ---. \textbf{Es la fuente de todas las cifras
  económicas del documento.}
  \item[profiles/] Evidencia de cartera de cada perfil, con la misma estructura.
\end{description}

\section{Los catorce controles de neutralización}

La neutralización de estilo del Capítulo~\ref{chap:resultados} regresa la señal contra estos
catorce controles antes de recalcular el Rank-IC: \texttt{factor\_pe}, \texttt{factor\_pb},
\texttt{factor\_ps}, \texttt{factor\_ev\_ebitda}, \texttt{factor\_relative\_return\_12m},
\texttt{factor\_momentum\_12\_1}, \texttt{factor\_realized\_vol\_63d},
\texttt{factor\_realized\_vol\_126d}, \texttt{factor\_beta\_252d}, \texttt{factor\_roe},
\texttt{factor\_roic}, \texttt{factor\_operating\_margin}, \texttt{factor\_eps\_growth\_yoy} y
\texttt{factor\_sales\_per\_share\_growth\_yoy}.

\section{Trazabilidad de las distribuciones nulas}

Los contrastes basados en simulación persisten tanto sus estadísticos-resumen como \emph{la
distribución que los generó}. Del test de permutación se conservan el $p$-valor, el Rank-IC
observado y los estadísticos de las 9.999 permutaciones; del \textit{bootstrap} por bloques, los
extremos del intervalo y las 2.000 medias remuestreadas; de las carteras aleatorias, los cuantiles
y el CAGR de cada una de las 1.000 simulaciones, tanto en la versión general como en la emparejada
por riesgo.

Esto permite que las Figuras~\ref{fig:permutacion},~\ref{fig:bootstrap} y~\ref{fig:aleatorias}
representen histogramas de la distribución nula y no solo sus resúmenes: se ve dónde cae el valor
observado dentro de la nube completa, no únicamente que la supera. Todo lo representado son valores
medidos y persistidos; ninguna figura simula una distribución plausible a partir de un resumen.

Los estadísticos-resumen se pueden recalcular exactamente desde las réplicas guardadas —el
$p$-valor del contraste de permutación y los extremos del intervalo \textit{bootstrap} se
reproducen sin discrepancia—, de modo que la figura y la cifra del texto proceden demostrablemente
del mismo cálculo.

Las réplicas se guardan únicamente en los tres contrastes de robustez que se representan. La
comparación pareada que decide cada variable del catálogo también usa \textit{bootstrap} por
bloques, pero ahí se descartan a propósito: se evalúa una vez por candidato de cada variable y
arrastrar sus remuestras multiplicaría el tamaño de \artefacto{decisions.json} sin que ninguna
figura las utilice.

\section{Alcance de las conclusiones}

Las conclusiones del texto principal se apoyan únicamente en las métricas y contrastes explicitados
en el Capítulo~\ref{chap:resultados}. El resto de los artefactos aquí inventariados se conserva para
auditoría y no sostiene ninguna afirmación que no aparezca respaldada en el cuerpo del documento.
